\clearpage
\begin{summarybox}
\textbf{\textcolor{CSummary}{一句话核心：对称割线张开，弦的方向不变。}}

过双曲线上 $P(u,v)$ 作斜率为 $m,-m$ 的两条真实割线，其另两个交点为 $M,N$。把 $P$ 关于 $x$ 轴翻成 $Q(u,-v)$：
\[
\boxed{MN\parallel Q\text{ 点切线}}
\]
若 $v\ne0$，$\displaystyle\boxed{k_{MN}=-\frac{b^2u}{a^2v}}$；若 $v=0$，弦和切线均竖直。

\medskip
\textbf{\textcolor{CSummary}{三秒识别流程}}
\[
\boxed{\text{同一双曲线点 }P}\ \longrightarrow\
\boxed{\text{两割线斜率 }m,-m}\ \longrightarrow\
\boxed{Q(u,-v)\text{ 的切线方向}}.
\]

\begin{itemize}[leftmargin=*,itemsep=0.42em]
 \item \textbf{先确定方向：} 直接写 $MN:\frac{ux}{a^2}+\frac{vy}{b^2}=t$，不要先求两交点坐标。
 \item \textbf{再确定位置：} 结合交点范围、内心所在区域、过点或距离条件找 $t$。
 \item \textbf{三个硬边界：} 两条割线\textbf{不同}，各有\textbf{另一个有限交点}，且 $v=0$ 时\textbf{不要套斜率分式}。
 \item \textbf{特别易错：} 两条割线对称，\textbf{不代表} $M,N$ 两个交点对称；$MN$ 与切线\textbf{平行而非一定重合}。
\end{itemize}

\medskip
\textbf{\textcolor{CSummary}{与其他结论的关系}}

P055 帮你看出“关于竖直线的反射使斜率变号”；C095 则把这个\textbf{直线方向信息}升级为\textbf{双曲线上弦与切线的联系}。使用时仍须核对公共点、第二交点及射线朝向。
\end{summarybox}
